\chapter{Pergunta, hipótese e objetivos}\label{pergunta-hipuxf3tese-e-objetivos}

\textbf{Pergunta principal:} em quais condições substituir a resposta
dependente da frequência por uma resposta avaliada em
\(f_{\mathrm{ref}}\) preserva a inferência da massa do gráviton, e
quando isso altera limites, evidências entre modelos ou a identificação
de uma componente escalar?

\textbf{Hipótese de trabalho:} como a dispersão depende de
\(m_g^2/f^2\), a adequação dessa substituição depende do espectro, da
janela temporal, dos pesos dos estimadores, do ruído e da proximidade do
limiar de propagação. Ela pode ser suficiente em alguns regimes e
produzir perda de informação ou viés em outros. A hipótese não pressupõe
o sinal ou o tamanho do efeito.

\textbf{Objetivo geral:} construir e validar uma comparação reproduzível
entre inferência com frequência explícita e inferência comprimida para
ondas gravitacionais massivas em PTAs.

\textbf{Objetivos específicos:}

\begin{enumerate}
\def\labelenumi{\arabic{enumi}.}
\tightlist
\item
  Reproduzir o núcleo cinemático do TG e explicitar sua relação com
  teorias massivas contemporâneas.
\item
  Implementar a resposta tensorial dispersiva com termos da Terra e do
  pulsar, controlando os limites analíticos.
\item
  Construir simulações nas quais o espectro, a massa, a geometria e o
  ruído sejam conhecidos.
\item
  Medir o efeito da compressão espectral, das escolhas a priori e da
  covariância entre estimadores.
\item
  Estender um subconjunto dos testes para tensor mais helicidade zero de
  Fierz--Pauli, respeitando a relação entre deformações escalar
  transversal e longitudinal.
\item
  Produzir um artigo sobre critérios de validade, informação
  efetivamente fornecida pelos dados e consequências para testes
  futuros.
\end{enumerate}
